\subsection{RQ0: How do \ToolTeralizer{} and JARVIS compare in applicability and value coverage?}
\label{sec:jarvis-comparison}

% P1: The two generalization mechanisms.
JARVIS~\cite{peleg_2018_jarvis} is the closest prior work to \ToolTeralizer{} (Section~\ref{sec:rw-generalization}).
It transforms conventional unit tests into property-based tests
by abstracting the constant values of repetitive tests,
selecting each abstraction from a manually defined template library
based on static analysis of test code alone~\cite{peleg_2018_jarvis}.
\ToolTeralizer{}, in contrast,
derives its input/output specifications
through static and dynamic analysis
of both the tests and the implementation of a target project,
and generalizes the developer-written assertions themselves (Section~\ref{sec:approach}).

% P2: Comparison basis, and the evaluation environment reconstruction.
The JARVIS implementation and abstraction template library are not publicly available,
so we benchmark \ToolTeralizer{} against the results reported in the JARVIS paper.
JARVIS was evaluated on the JUnit test suites
of 12 Apache Commons projects~\cite{peleg_2018_jarvis}.
However, the JARVIS paper identifies neither the exact project revisions it evaluated
nor the exact date of its evaluation.
In a best-effort attempt to reconstruct the evaluation environment,
we therefore pin each of the 12 projects
to its last commit before February~1, 2017 ---
the date on which the commons-lang fix that JARVIS reports was submitted.
We generalize the test suites of all 12 projects
with the \VariantImproved{} strategy at 100 \tries{},
matching the 100 samples per synthesized property used during the original evaluation of JARVIS.

% P3: Why the two outputs differ in unit.
% Because the two approaches follow different generalization mechanisms,
% their results must be aligned on a shared unit before they can be compared.
% JARVIS groups structurally repetitive tests into scenarios,
% that is, tests that execute the same statement sequence
% and differ only in their constant values,
% and it reports its results per scenario.
% \ToolTeralizer{} creates one PBT per generalizable assertion
% (Section~\ref{sec:generalized-test-creation}).
% Thus, JARVIS's scenario-based PBTs
% and \ToolTeralizer{}'s assertion-based PBTs
% are not directly comparable.

% P4: Why the two sets of results are aligned at the method under test.
% What both approaches share is the set of methods whose behaviors they exercise.
% Every test in a scenario targeted by JARVIS executes the same tested method calls,
% and \ToolTeralizer{} maps every assertion it targets to exactly one tested method call.
% Thus, we aggregate generalization results at the level of these methods under test (MUTs).
% The \TzRqzeroTabletwoReportedRows{} scenarios that JARVIS reports
% correspond to \TzRqzeroTabletwoDistinctMuts{} distinct MUTs (Table~\ref{tab:jarvis-scenario-pvc}):
% the 3 PolynomialFunction scenarios all call \texttt{PolynomialFunction.value},
% \texttt{testMinMaxDouble} calls both \texttt{FastMath.min} and \texttt{FastMath.max},
% and each of the remaining 6 scenarios calls one method.

% P5: Roadmap.
The comparison proceeds in three parts:
the first compares parameter value coverage
of JARVIS and \ToolTeralizer{} for the \TzRqzeroTabletwoReportedRows{} scenarios
for which JARVIS successfully produces property-based tests,
the second compares applicability
across all 12 projects targeted in the JARVIS evaluation
by measuring how many MUTs each tool produces PBTs for,
and the third examines the relationship
between parameter value coverage and test suite effectiveness.

\subsubsection{Parameter Value Coverage on the Scenarios that JARVIS Successfully Generalizes}%
\label{sec:jarvis-scenario-pvc}

% P1: Measurement and table contract.
Table~\ref{tab:jarvis-scenario-pvc} compares parameter value coverage (PVC):
the number of distinct input values a test suite exercises
for the parameters of a MUT~\cite{sampath_2008_prioritizing}.
PVC is the primary coverage metric reported by the JARVIS evaluation~\cite{peleg_2018_jarvis}.
The last three columns show PVC before generalization,
after generalization with JARVIS,
and after generalization with \ToolTeralizer{}.
Each row represents one of the \TzRqzeroTabletwoReportedRows{} scenarios
successfully generalized by JARVIS.

% Data source: data/generated/rq0/rq0-table2-comparison.csv
\begin{table}[tb]
  \caption{Parameter value coverage (PVC) before generalization, after generalization with JARVIS, and after generalization with \ToolTeralizer{} across the 10 scenarios reported in the JARVIS evaluation.}
  \label{tab:jarvis-scenario-pvc}
  \centering
    \begin{tabular*}{\linewidth}{@{\hspace{\tabcolsep}\extracolsep{\fill}}rlrrr}
  \toprule
   &  & Original & JARVIS & \ToolTeralizer{} \\
  \# & JARVIS scenario & CUT PVC & PBT PVC & PBT PVC \\
  \midrule
  1 & \texttt{isAscii} & 6 & 59 & 230 \\
  2 & \texttt{isPrintable} & 195 & 45 & 197 \\
  3 & \texttt{testMinMaxDouble} & 9 & 400 & 368 \\
  4 & \texttt{toIntExact} & 2,001 & 65 & 2,074 \\
  5 & \texttt{IntervalTest} & 2 & 2 & 172 \\
  6 & \texttt{testConstants} & 5 & 105 & -- \\
  7 & \texttt{testfirstDerivativeComparison} & 7 & 264 & 105 \\
  8 & \texttt{testLinear} & 5 & 160 & 104 \\
  9 & \texttt{PrecisionTest} & 8 & 102 & -- \\
  10 & \texttt{testAbs} & 5 & 506 & -- \\
  \bottomrule
  \end{tabular*}

\end{table}


% P2: Head-to-head result and comparable sampling budget.
\ToolTeralizer{} creates generalized tests for
\TzRqzeroTabletwoSoundTabletwoRows{} of the
\TzRqzeroTabletwoReportedRows{} scenarios.
For 4 of these \TzRqzeroTabletwoReportedRows{} scenarios (1, 2, 4, and 5),
the PVC measured for \ToolTeralizer{} is higher than the PVC reported by JARVIS.
For the other 3 scenarios (3, 7, and 8), its PVC is lower.
Both evaluations exercise each synthesized PBT with 100 inputs.
The per-PBT sampling budget is therefore the same,
but the total budget for a scenario depends on how many PBTs it yields.
JARVIS may create multiple PBTs to preserve distinct constraints
or positive and negative regions found in its original tests.
\ToolTeralizer{} instead creates one PBT per generalizable assertion.

% P3: Cases where JARVIS reports higher PVC.
The number of PBTs created by the two tools
for each scenario explains some of the JARVIS advantages.
Scenarios 7 and 8 each have one numeric parameter.
Their JARVIS PVCs of 264 and 160 therefore require at least 3 and 2
executions of a 100-input PBT, respectively.
However, \ToolTeralizer{} only produces one PBT for each of these scenarios,
because each one only offers a single generalizable assertion.
Scenario 3 shows that the number of created PBTs
is not the only factor that affects PVC results.
\ToolTeralizer{} produces 2 PBTs in this case:
one for \texttt{FastMath.min} and one for \texttt{FastMath.max}.
However, the input samples generated by these PBTs repeat some values.
Those repetitions reduce PVC independently of PBT count.
Furthermore, JARVIS's inability to sample \texttt{Double.NaN}
also reduces its coverage for this scenario~\cite{peleg_2018_jarvis},
contributing yet another modulating factor.

% P4: Cases where Teralizer reports higher PVC.
The higher \ToolTeralizer{} values primarily reflect input retention.
By design, \ToolTeralizer{}'s generalized test suites retain the original input values
when sampling input values for its PBTs
(Section~\ref{sec:three-variant-design}).
JARVIS constrains its generators to regions that contain the original examples,
but random sampling does not guarantee that these examples are executed again.
This distinction explains Scenarios 2 and 4.
JARVIS reduces their PVC from 195 to 45 and from 2,001 to 65,
whereas \ToolTeralizer{} retains the original values
and increases the corresponding PVC to 197 and 2,074, respectively.
For Scenario 1, JARVIS increases PVC from 6 to 59,
yet \ToolTeralizer{} achieves a larger increase to a PVC of 230
by creating multiple PBTs rather than just a single one
--- one for each assertion in the original test.

% P5: Truncated JARVIS run and interpretation boundary.
% Scenario 5 is a separate case.
% Its JARVIS PVC of 2 records a PBT that terminates
% on the second generated input on average because the property fails.
% It is therefore a truncated run rather than ordinary 100-input coverage~\cite{peleg_2018_jarvis}.
% Across the 7 shared scenarios, the observed differences follow 4 mechanisms:
% PBT count, retained original inputs, repeated sampled values,
% and premature property termination.
% Exact JARVIS PBT counts are not reported,
% but the published PVC values establish the lower bounds above.

% P6: Why Teralizer declines 3 scenarios.
\ToolTeralizer{} does not create generalized tests for the 3 remaining scenarios (6, 9, and 10)
because specification extraction does not establish a usable relation
between the generalized input and the expected output for these cases.
In Scenarios 9 and 10, the corresponding implementations
of \texttt{Precision.equals} and \texttt{Abs.value}
compute their outputs through raw IEEE-754 bit operations.
\ToolSPF{} models floating-point values as real numbers
and thus cannot preserve the symbolic input/output relation through these operations.
In Scenario 6, the method's output depends on internal state
rather than its input parameter.
The extracted specification tracks the input parameter but not this internal state,
so it cannot establish the input/output relation needed for a generalized oracle.
RQ5 and RQ6 (Sections~\ref{sec:limitations-eval} and~\ref{sec:limitations-eval-extended})
analyze the prevalence of the different exclusion mechanisms in further detail.

\subsubsection{Applicability Across the Full JARVIS Benchmark}%
\label{sec:jarvis-applicability}

% P1: Population and reported applicability.
The \TzRqzeroTabletwoReportedRows{} scenarios reported in Table~\ref{tab:jarvis-scenario-pvc}
are all found within 2 of the 12 Apache Commons projects
used during the original evaluation of JARVIS:
commons-lang and commons-math~\cite{peleg_2018_jarvis}.
For the 10 remaining projects that are part of the JARVIS benchmark,
the paper reports no successful generalization.
Table~\ref{tab:jarvis-applicability} compares this reported applicability result
with \ToolTeralizer{}'s measured results across the same 12 projects.
A dash in the JARVIS columns means
that the paper reports no PBT for the project.
A dash in the \ToolTeralizer{} columns means
that the 12-project evaluation created no generalized test for the project.

% Data source: data/generated/rq0/rq0-breadth-summary.csv
\begin{table}
  \caption{MUTs with a generalized test and the PVC of those tests, per project.}
  \label{tab:jarvis-applicability}
  \centering
    \begin{tabular*}{\linewidth}{@{\hspace{\tabcolsep}\extracolsep{\fill}}lrrrr}
  \toprule
   & \multicolumn{2}{r}{JARVIS} & \multicolumn{2}{r}{\ToolTeralizer{}} \\
  \cmidrule(lr){2-3} \cmidrule(lr){4-5}
  Benchmark project & PBT PVC & MUTs & PBT PVC & MUTs \\
  \midrule
  commons-math & 1,604 & 7 & 12,471 & 90 \\
  commons-lang & 104 & 2 & 26,079 & 208 \\
  commons-cli & -- & -- & 217 & 2 \\
  commons-codec & -- & -- & 983 & 6 \\
  commons-collections & -- & -- & 102 & 2 \\
  commons-configuration & -- & -- & 1,403 & 9 \\
  commons-csv & -- & -- & 4 & 1 \\
  commons-email & -- & -- & -- & -- \\
  commons-io & -- & -- & 774 & 7 \\
  commons-jexl & -- & -- & -- & -- \\
  commons-pool & -- & -- & -- & -- \\
  commons-text & -- & -- & 1,144 & 12 \\
  \midrule
  all 12 projects & 1,708 & 9 & 43,177 & 337 \\
  \bottomrule
  \end{tabular*}

\end{table}


% P2: The unit that carries both sets of results.
% The two sets of results do not share a unit.
% JARVIS reports per scenario, its grouping of structurally repetitive tests,
% and the paper publishes no count
% of the property-based tests it creates.
% \ToolTeralizer{} instead creates one PBT per generalizable assertion
% (Section~\ref{sec:generalized-test-creation}).
% Both tools do identify the methods whose behavior they exercise,
% so Table~\ref{tab:jarvis-applicability} counts MUTs.
% The \TzRqzeroTabletwoReportedRows{} JARVIS scenarios call
% \TzRqzeroTabletwoDistinctMuts{} of them:
% one scenario calls both \texttt{FastMath.min} and \texttt{FastMath.max},
% whereas 3 scenarios share \texttt{PolynomialFunction.value}.

% P3: Applicability result.
\ToolTeralizer{} creates generalized tests for
\TzRqzeroBreadthTeralizerTotalSoundMuts{} MUTs in 9 of the 12 projects,
whereas JARVIS reports successful generalization for only
\TzRqzeroBreadthJarvisTotalMuts{} MUTs
in \TzRqzeroBreadthPublishedProjects{} of the 12 projects.
% P4: What JARVIS requires before generalization can succeed.
Three requirements that are inherent to the JARVIS approach limit its applicability.
First, abstraction requires at least two structurally matching tests.
As reported in the JARVIS paper, the median project-level share of scenarios
that meet this repetition prerequisite is 30.6\% across the 12 projects,
with a range of 13.3--60.0\%~\cite{peleg_2018_jarvis}.
Second, JARVIS selects each abstraction
from a manually curated template library,
so a fitting template must exist for the relation between the observed inputs.
Third, JARVIS uses Z3 to instantiate the coefficients of the selected template.
Consequently, solver theory limitations (Section~\ref{sec:symbolic-analysis})
affect which data types the template library can usefully cover.

% P5: How Teralizer avoids these three requirements.
\ToolTeralizer{} avoids these three requirements
by deriving its specifications via single-path symbolic analysis
rather than static analysis of repetitive tests.
A single test execution provides both the path condition
and the symbolic output expression,
so \ToolTeralizer{} needs neither a second matching test
nor a curated template library.
\ToolTeralizer{} also performs no constraint solving at any pipeline stage,
so the types it can generalize depend on \ToolSPF{}'s symbolic representations
instead of solver theories.
Section~\ref{sec:when-it-fails} examines why this distinction matters in practice.

% P6: The 3 projects without generalized tests.
In the 3 projects without generalized tests,
\ToolTeralizer{} fails to obtain a specification that it can generalize.
Across commons-email and commons-jexl,
filtering and \ToolSPF{} failures exclude every assertion.
In commons-email, most failures occur while \ToolSPF{} initializes \texttt{java.net.IDN},
which the tested code uses to convert internationalized addresses.
Its initializer reads a Unicode data file from the class path,
which \ToolSPF{} does not provide, so the initializer fails its own assertion.
In commons-jexl, most failures occur because the test constructor
requests the default charset,
which the JDK derives from the \texttt{file.encoding} property.
\ToolSPF{} leaves that property unset,
so the request fails before the MUT is called.
For commons-pool, the pipeline stops before specification extraction,
because the original test suite does not finish
within the 300-second execution budget.
RQ5 and RQ6 analyze these failure causes in detail
(Sections~\ref{sec:limitations-eval} and~\ref{sec:limitations-eval-extended}).

\subsubsection{Interpreting Parameter Value Coverage}%
\label{sec:jarvis-pvc-budget}

% P1: Why PVC cannot measure the effectiveness that the evaluation targets.
A comprehensive evaluation of test generalization approaches
should address both applicability
and regression detection effectiveness.
The preceding section measures applicability
by counting the number of projects and MUTs
for which JARVIS and \ToolTeralizer{} create generalized tests.
However, PVC --- which is the primary metric used in the JARVIS paper ---
measures input diversity rather than fault detection effectiveness:
it counts distinct input values without observing whether they expose faulty behavior
\cite{sampath_2008_prioritizing,inozemtseva_2014_coverage,peleg_2018_jarvis}.

% P2: Construction of the budget experiment.
We examine the distinction between input diversity and fault detection effectiveness
for the 10 scenarios of the initial PVC evaluation (Table~\ref{tab:jarvis-scenario-pvc}).
For this, we reuse the generalized tests created by \ToolTeralizer{}'s
\VariantImproved{} generalization strategy
and increase their sampling budget from 100 to 200 and 1,000 \tries{}.
This setup raises PVC by exercising more distinct values
without changing the test code or how inputs are selected.
We then use covered mutation score ---
which is commonly accepted as a reasonable proxy
for fault detection capability (Section~\ref{sec:mutation-testing})
--- to observe whether the additional samples detect a larger number of seeded faults.
Table~\ref{tab:jarvis-pvc-budget} reports the mutation results at each budget.

% Data source: data/generated/rq0/rq0-pvc-budget.csv
\begin{table}
  \caption{PVC and mutation testing results for the same 10 generalized tests at 3 sampling budgets. Only the sampling budget changes between the rows. PVC sums the distinct input values that the tests exercise, and Score is the killed mutants divided by the covered mutants.}
  \label{tab:jarvis-pvc-budget}
  \centering
    \begin{tabular*}{\linewidth}{@{\hspace{\tabcolsep}\extracolsep{\fill}}lrrrrr}
  \toprule
  \VariantImproved{} &  &  & \multicolumn{3}{r}{Mutation Testing Results} \\
  \cmidrule(lr){4-6}
  Sampling Budget & Tests & PVC & Killed & Covered & Score \\
  \midrule
  100 tries & 10 & 1,135 & 51 & 80 & 63.8\% \\
  200 tries & 10 & 2,152 & 51 & 80 & 63.8\% \\
  1,000 tries & 10 & 10,055 & 51 & 80 & 63.8\% \\
  \bottomrule
  \end{tabular*}

\end{table}


% P3: Observed contrast and the handoff to the remaining evaluation.
At both the 200 and 1,000 \tries{} budgets, PVC exceeds its 100-try baseline.
Nevertheless, all 3 budgets cover the same
\TzRqzeroBudgetImprovedOnezerozeroTriesCoveredMutants{} mutants
and kill the same
\TzRqzeroBudgetImprovedOnezerozeroTriesKilledMutants{}.
They therefore achieve the same covered mutation score of
\TzRqzeroBudgetImprovedOnezerozeroTriesCoveredMutationScore{}.
Thus, while PVC remains useful for describing input diversity
and for the direct comparison with JARVIS
--- which lacks a publicly available implementation and thus
prevents a more robust effectiveness comparison ---
the remaining effectiveness analyses use mutation score instead,
which has stronger supporting evidence as a regression detection proxy
(Section~\ref{sec:mutation-testing}).

\rqanswerbox{0}{
  Across the 12 projects of the JARVIS benchmark,
  \ToolTeralizer{} creates generalized tests
  for \TzRqzeroBreadthTeralizerTotalSoundMuts{} MUTs in 9 projects,
  whereas JARVIS reports successful generalization for
  only \TzRqzeroBreadthJarvisTotalMuts{} MUTs
  in \TzRqzeroBreadthPublishedProjects{} projects.
  Among the 10 scenarios that JARVIS successfully generalizes,
  \ToolTeralizer{} creates PBTs with higher PVC in 4,
  lower PVC in 3, and fails to create PBTs in 3 scenarios.
  This suggests that \ToolTeralizer{} is more broadly applicable than JARVIS
  while neither approach has a clear input diversity advantage.
  However, it does not provide definitive evidence
  about the fault detection effectiveness of either approach
  because PVC is a comparatively weak predictor for this.
}
